\documentclass[11pt,a4paper]{article}
\usepackage[margin=2.5cm]{geometry}
\usepackage{hyperref}
\usepackage{microtype}
\usepackage{parskip}

\begin{document}

\noindent
To the Editor,\\[1em]

We submit for your consideration ``From Microprice to Microdiffusion: Heavy-Tailed
Price Diffusion in Limit Order Books.''

\bigskip
\noindent\textbf{What the paper does.}
Standard microstructure models condition the price \emph{drift} on order-book state,
most directly through the microprice. We ask the complementary question: does the
order-book state also condition the \emph{second moment} of price changes? We show that
it does. Using tick-by-tick data from the Qatar Stock Exchange (36 instruments, 216
trading days), we estimate a conditional diffusion surface $a_{xx}(I,S)$ indexed by
order-book imbalance $I$ and relative spread $S$, and study whether and how it
transfers out of sample.

\bigskip
\noindent\textbf{Main findings.}
\begin{enumerate}

\item \emph{The surface is stable and portable.} Trained on one set of
  days and evaluated on disjoint days, the surface transfers with a Spearman rank
  correlation of $0.48$ (instrument-clustered 95\% interval $[0.40,\,0.56]$). The
  spread axis carries most of the transferable level; the within-spread imbalance
  shape is a real contemporaneous effect but does not reliably transfer day to day.
  The surface also transfers across volatility regimes: cross-regime rank correlations
  of $0.849$ and $0.845$ exceed the adjacent-block baseline of $0.785$.

\item \emph{A Generalized Hyperbolic innovation is needed.} The Gaussian collapses in
  the far tail by orders of magnitude; the maximum-likelihood Student-$t$ is pulled to
  $\hat{\nu}=1.56$ (no finite variance). The Generalized Hyperbolic matches the
  empirical tail closely and improves an intraday VaR backtest materially over the
  Gaussian: at $\alpha=0.01$ the median violation rate falls from $0.0275$ to
  $0.0133$ across 36 instruments.

\item \emph{The model is competitive with time-series benchmarks without using return
  history.} Against jointly estimated GARCH-$t$, the model is a statistical tie rather
  than a clear winner. Its value is different: it gives an observable and transferable
  book-state surface that can be read before the next move occurs.

\end{enumerate}

\bigskip
\noindent\textbf{Fit for this journal.}
The paper is written for readers interested in market microstructure, financial
econometrics, and short-horizon risk measurement. It introduces a new estimand, the
conditional diffusion surface, shows that it transfers across days, instruments, and
volatility regimes, and uses it to build a heavy-tailed predictive distribution for
short-horizon price risk. We are submitting it to [Journal name] because the journal
regularly publishes empirical work on forecasting, risk measurement, and financial
market dynamics.

\bigskip
\noindent Geoffrey Ducournau, Yibo Wang, Jinliang Li\\
\noindent jll@tsinghua.edu.cn

\end{document}
